\chapter{Аудит кандидатов кварк--лептонного коннектора
\texorpdfstring{$H_{15}$}{H15}}
\label{ch:v10-h15-quark-lepton-connector-audit}

Предыдущий гейт показал, что абстрактное ребро $Q_L-L_L$ делает типовой
граф связным. Теперь возвращаем квантовые числа. Для порядка вершин
$(Q_L,L_L,u_R,d_R,e_R)$ положим
\begin{equation}
 b_{QL}=(1,-1,0,0,0)^T.
\end{equation}
Оно действительно повышает ранг incidence с $3$ до $4$, но соединяет две
левые вершины. Для знака хиральности $\chi=(1,1,-1,-1,-1)$ имеем
\begin{equation}
 \chi b_{QL}=0.
\end{equation}
Кроме того, цветовая метка даёт дефект $1$, а целочисленная метка $6Y$ ---
дефект $4$. Поэтому это полезное графовое ребро не является допустимым
блоком нечётного калибровочно-инвариантного оператора Дирака.

Минимальное физически типизированное завершение обязано идти между
противоположными хиральностями. Полный набор отсутствующих рёбер равен
\begin{equation}
 \{L_L-u_R,\ L_L-d_R,\ Q_L-e_R\}.
\end{equation}
Единственная пара, читаемая как одно комплексное поле, есть
\begin{equation}
 R_2\sim(\mathbf3,\mathbf2)_{7/6},\qquad
 \{L_L-u_R,\ Q_L-e_R\}.
\end{equation}
Она делает граф связным и создаёт ровно один смешанный цикл.

Однако бимодульные координаты
\begin{equation}
 Q_L=(\mathbb H,M_3),\quad L_L=(\mathbb H,\mathbb C),\quad
 u_R,d_R=(\mathbb C,M_3),\quad e_R=(\mathbb C,\mathbb C)
\end{equation}
показывают строгий запрет. Оба $R_2$-ребра меняют сразу левую и правую
координаты. Тем самым они нарушают условие первого порядка на неизменённой
геометрии Стандартной модели. Real-сопряжение меняет координаты местами, но
не устраняет двойной дефект.

Аудит одиннадцати маршрутов по шести независимым критериям имеет ранг $6$.
Ни один маршрут не набирает $6/6$. Унаследованный Higgs-лес набирает $5/6$,
но не связывает кварковую и лептонную компоненты. Target-loaded обобщённый
$R_2$ также набирает $5/6$, но не имеет унаследованного носителя. Обычный
$R_2$ набирает $3/6$: типизация и связность верны, строгий первый порядок
и происхождение отсутствуют.

\begin{theorem}[Разделение графового и физического коннектора]
Одно ребро $Q_L-L_L$ достаточно лишь для связности типового графа, но
невозможно как нечётный калибровочно-инвариантный Dirac-блок. Минимальное
типизированное завершение есть двухрёберный $R_2$-прямоугольник; на
фиксированной алгебре, представлении и строгом условии первого порядка оба
его ребра запрещены.
\end{theorem}

\begin{nogocriterion}[Связность не заменяет бимодульную допустимость]
Нельзя использовать условное $Q_L-L_L$-ребро предыдущего гейта как
физический источник uniform amplification. Нельзя также объявить $R_2$
унаследованным, пока явно не построено обобщённое условие первого порядка
или расширение конечной геометрии и независимая нормировка коннектора.
\end{nogocriterion}

\begin{protocol}\begin{enumerate}
\item прямое $Q_L-L_L$-ребро связывает граф: \textbf{да};
\item оно имеет правильную нечётную хиральность: \textbf{нет};
\item его цветовой и $6Y$-дефекты: \textbf{$1$ и $4$};
\item минимальное типизированное завершение: \textbf{$R_2$, два ребра};
\item смешанных циклов после завершения: \textbf{$1$};
\item оба $R_2$-ребра проходят строгий первый порядок: \textbf{нет};
\item прошедших кандидатов: \textbf{$0/11$};
\item physical origin: \textbf{$0/3$};
\item следующий гейт: \textbf{обобщённый первопорядковый родитель $R_2$}.
\end{enumerate}\end{protocol}

Машинный сертификат сохранён в
\path{s2t_v10_particle_wrinkle_dislocation_callias_equal_charge_twist_m4_h15_quark_lepton_connector_candidate_audit_gate_results.json}.